\chapter{Theory and Foundations}
\begin{definition}[Graph]
A finite graph is a tuple
\[
G=(V,E),
\]
where \(V=\{v_1,\dots,v_n\}\) is a finite set of vertices. For an undirected
graph, \(E\) is a set of unordered pairs \(\{u,v\}\) with \(u,v\in V\). For a
directed graph, \(E\subseteq V\times V\) is a set of ordered pairs
\((u,v)\).
\end{definition}

\begin{definition}[Attributed graph]
An attributed graph is a graph \(G=(V,E)\) equipped with node and (optionally) edge attributes.
We denote node attributes by a function
\[
\mathbf{x}: V \to \mathbb{R}^{d_x}, \qquad v \mapsto \mathbf{x}_v,
\]
and edge attributes by a function
\[
\mathbf{e}: E \to \mathbb{R}^{d_e}, \qquad (u,v) \mapsto \mathbf{e}_{uv}.
\]
For categorical attributes, \(\mathbf{x}_v\) and/or \(\mathbf{e}_{uv}\) may take values in a finite set (e.g., atom/bond types).
\end{definition}

\begin{definition}[Adjacency matrix]
Let \(G=(V,E)\) be a graph with \(|V|=n\) and an ordering of vertices.
The adjacency matrix \(A \in \mathbb{R}^{n \times n}\) is defined by
\[
A_{ij} =
\begin{cases}
1, & \text{if } \{v_i,v_j\}\in E \text{ for an undirected graph},\\
1, & \text{if } (v_i,v_j)\in E \text{ for a directed graph},\\
0, & \text{otherwise.}
\end{cases}
\]
For weighted graphs, \(A_{ij}\) may store an edge weight \(w_{ij}\) instead of a binary indicator.
\end{definition}

\begin{definition}[Degree matrix]
The degree matrix \(D \in \mathbb{R}^{n \times n}\) is the diagonal matrix with entries
\[
D_{ii} = \sum_{j=1}^n A_{ij}, \qquad D_{ij}=0 \text{ for } i \neq j.
\]
For directed graphs, one may analogously define in-degree and out-degree matrices.
\end{definition}

\begin{definition}[Graph Laplacian]
The (unnormalised) graph Laplacian is
\[
L = D - A.
\]
A commonly used normalized variant is the symmetric normalised Laplacian
\[
L_{\mathrm{sym}} = I - D^{-1/2} A D^{-1/2},
\]
defined when \(D\) is invertible on the support of the graph.
\end{definition}
